\clearpage
\begingroup
\fancyhead[L]{\small THEORIE DER BEWEGUNG DES MONDES.}
\renewcommand{\headrulewidth}{0pt}
\SourceLeaf{623}
\section*{ZWEITER ABSCHNITT.}
\begin{center}
\itshape Integration der Fundamentalgleichungen, mit Beiseitesetzung der
Störungskräfte der Sonne.
\end{center}

\par\noindent\hspace*{\parindent}\textls[120]{Lehrsatz aus der Integralrechnung:}

\par\noindent\hspace*{\parindent}Das Integral der Differentialgleichung des zweiten Grades
\[
0=\frac{\mathrm{dd}w}{\mathrm{d}v^2}+\alpha\alpha w+Z,
\]
\noindent wo \(Z\) eine Function von \(v\) und \(\alpha\alpha\) eine beständige
positive Grösse bedeutet, ist
\[
\begin{aligned}
\alpha w&=\cos\alpha v\int\sin\alpha v\mathpunct{.}\,Z\mathrm{d}v
 -\sin\alpha v\int\cos\alpha v\mathpunct{.}\,Z\mathrm{d}v\\
 &\qquad+\mathit{Const.}\sin(\alpha v+\mathit{Const.}).
\end{aligned}
\]

\par\noindent\hspace*{\parindent}Ist also \(Z=a+b\sin\beta v+c\sin\gamma v+\mathrm{etc.}\), so wird
\[
\begin{aligned}
w&=-\frac{a}{\alpha\alpha}+\frac{b}{\beta\beta-\alpha\alpha}\sin\beta v
 +\frac{c}{\gamma\gamma-\alpha\alpha}\sin\gamma v\\
 &\qquad+\mathrm{etc.}+\mathit{Const.}\sin(\alpha v+C.).
\end{aligned}
\]

\par\noindent Oder ist \(Z=a+b\cos\beta v+c\cos\gamma v+\mathrm{etc.}\), so wird
\[
\begin{aligned}
w&=-\frac{a}{\alpha\alpha}+\frac{b}{\beta\beta-\alpha\alpha}\cos\beta v
 +\frac{c}{\gamma\gamma-\alpha\alpha}\cos\gamma v\\
 &\qquad+\mathrm{etc.}+\mathit{Const.}\cos(\alpha v+C.).
\end{aligned}
\]

\par\noindent\hspace*{\parindent}Enthielte hingegen \(Z\) einen Theil \(a\cos\alpha v\), so würde dafür in
\(w\) zu setzen sein: \(-\tfrac12\frac{av\sin\alpha v}{\alpha}
-\tfrac14\frac{a\cos\alpha v}{\alpha\alpha}\), oder für einen Theil
\(a\sin\alpha v\) in \(Z\), käme in \(w\):
\(\tfrac12\frac{av\cos\alpha v}{\alpha}
-\tfrac14\frac{a\sin\alpha v}{\alpha\alpha}\mathpunct{\raisebox{.4ex}{.}}\)
Weiss man also a priori, dass \(w\) keine Cirkelbogen enthalten kann, so darf
auch \(Z\) keinen Sinus oder Cosinus von \(\alpha v\) enthalten.

\par\noindent\hspace*{\parindent}Man bemerkt bei obigen Fundamentalgleichungen, dass jede aus zwei
Theilen besteht, wovon der eine die Grössen \(\pi,\omega,\psi\) als Factor
enthält, der andere von ihnen unabhängig ist. Würde die Bewegung des Mondes
nicht von der Sonne gestört, so würden jene ersten Theile wegfallen, und die
Inte\SourceLeaf{624}gration\marginpar{\scriptsize VII.\quad78} sodann keine Schwierigkeit haben. Auch sieht man nach obigem
Lehrsatze, dass selbst mit Beibehaltung jener ersten Glieder die Integration
von Statten gehen würde, wenn man dieselben durch Functionen von \(v\)
darstellen könnte; allein dies geht nicht an, ohne die Bewegung des Mondes
schon zu kennen. Glücklicher Weise aber sind die Wirkungen der Sonne nicht
sehr beträchtlich; die grösste periodische Gleichung für die Länge des Mondes
beläuft sich nur auf \(1\tfrac13\) Grad, und das Apogäum des Mondes, welches
ohne die Einwirkung der Sonne unverrückt bleiben würde, erhält durch diese
nur eine langsame progressive Bewegung, die sich zur mittlern Bewegung des
Mondes wie 1 zu 118 verhält; die rücklaufende Bewegung des Knoten ist noch
unbeträchtlicher. Diese Umstände, ohne welche die Berechnung der Bewegung
des Mondes die Kräfte der Analyse überschreiten würde, machen es möglich,
sich der wahren Bewegung des Mondes stufenweise zu nähern, indem man
zuvörderst in den von \(\pi,\omega,\psi\) abhängigen Theilen die blosse
elliptische Bewegung des Mondes substituirt, und daraus eine erste genäherte
Bewegung ableitet; diese von neuem in jenen Theilen anstatt der elliptischen
gebraucht, um eine genauere Bestimmung zu erhalten, und diese successiven
Verbesserungen so lange wiederholt, als der Grad der Genauigkeit, den man
sich vorsetzt, erfordert.

\par\noindent\hspace*{\parindent}Nach diesen vorläufigen Bemerkungen nehmen wir also zuerst obige
Gleichungen mit Beiseitesetzung der Grössen \(\pi,\omega,\psi\) vor, um die
elliptische Bewegung zu erhalten, wobei es auch erlaubt ist, solche Grössen,
als von höhern Potenzen der Excentricität und Inclination abhängen, zu
vernachlässigen.

\par\noindent\hspace*{\parindent}Die dritte Gleichung wird hier
\[
\frac{\mathrm{dd}\theta}{\mathrm{d}v^2}+\theta=0,
\]
\noindent wovon das Integral (genau)
\[
\theta=g\sin\lambda,\quad\text{so dass}\quad\lambda=v+\mathit{Const.}
\]
\noindent Dieses zeigt, dass die Bahn des Mondes, vom Mittelpunkt der Erde
gesehen, ein grösster Kreis [ist], dessen Neigung gegen die Ekliptik zur
Tangente \(g\) hat; die Länge des Knoten \(=-\mathit{Const.}\) oder \(\lambda\)
das Argument der Breite\HistoricalNoteEight{Aber nicht in der Bahn des Mondes,
sondern auf der Ekliptik gezählt.]}.

\SourceLeaf{625}
\par\noindent\hspace*{\parindent}Die zweite Gleichung wird
\[
0=\frac{\mathrm{dd}u}{\mathrm{d}v\mathrm{d}u}
 +\frac{2\mathrm{d}p}{p\mathrm{d}v},
\]
\noindent woraus
\[
\log\frac{pp\mathrm{d}u}{\mathrm{d}v}=\log\mathit{Const.},
\]
\noindent also
\[
\frac{pp\mathrm{d}u}{\mathrm{d}v}=\mathit{Const.}=A.
\]
\noindent Hienach wird die erste Gleichung
\[
0=\frac{\mathrm{dd}p}{\mathrm{d}v^2}+p-AAk(1+\theta\theta)^{-\frac32},
\]
\noindent wenn man Kürze halber \(\frac{(m+m')\mathrm{d}t^2}{\mathrm{d}u^2}=k\)
setzt; folglich, obigen Werth von \(\theta\) substituirt, und die höhern
Potenzen von \(g\) vernachlässigt:
\[
0=\frac{\mathrm{dd}p}{\mathrm{d}v^2}+p
-AAk\left(1-\tfrac34gg+\tfrac34gg\cos2\lambda\right).
\]
\noindent Hieraus \(p=AAk(1-\tfrac34gg)(1-e\cos M-\tfrac14gg\cos2\lambda)\).
Also, den beständigen Theil von \(p\) oder \(AAk(1-\tfrac34gg)=h\) gesetzt,
\[
\begin{aligned}
p&=h(1-e\cos M-\tfrac14gg\cos2\lambda)\\[5pt]
\log p&=\log h-\tfrac14ee-e\cos M-\tfrac14ee\cos2M-\tfrac14gg\cos2\lambda\\[5pt]
\log pp\frac{\mathrm{d}u}{\mathrm{d}v}
 &=\log A=\tfrac12\log h-\tfrac12\log k+\tfrac38gg\\[5pt]
\log\frac{\mathrm{d}u}{\mathrm{d}v}
 &=-\tfrac32\log h-\tfrac12\log k+\tfrac12ee+\tfrac38gg\\
 &\qquad+2e\cos M+\tfrac12ee\cos2M+\tfrac12gg\cos2\lambda\\[5pt]
\frac{\mathrm{d}u}{\mathrm{d}v}
 &=h^{-\frac32}k^{-\frac12}(1+\tfrac32ee+\tfrac38gg)\\
 &\qquad(1+2e\cos M+\tfrac32ee\cos2M+\tfrac12gg\cos2\lambda),
\end{aligned}
\]
\noindent folglich
\[
1+ee+\tfrac14gg=hk^{\frac13}
\]
\noindent und
\[
u=v+2e\sin M+\tfrac34ee\sin2M+\tfrac14gg\sin2\lambda,
\]
\noindent wo man alle Grössen vernachlässigt hat, die höhere Potenzen von
\(e\) und \(g\) enthalten. Übrigens ist hier offenbar \(e\) eine durch die
Beobachtungen zu bestimmende Constante (die Excentricität), und \(M\) die
wahre Anomalie\HistoricalNoteEight{Aber nicht in der Bahn des Mondes,
sondern auf der Ekliptik gezählt.]} \(=v-\mathit{Const.}\), so dass
\(\mathit{Const.}\) die unveränderliche Länge der Erdferne [ist].

\par\medskip\noindent\hfill\rule{.19\textwidth}{.3pt}\hfill\mbox{}\par
\noindent\hfill 78\textsuperscript{*}\par
\clearpage
\endgroup
